% Table 4: 消融实验——各组件的作用
\begin{table}[t]
\centering
\caption{Ablation study: contribution of each SESS component (Qwen3-4B, WildJailbreak). The Skill System (CS + retrieval + maintenance) provides the primary gain (+22.7pp). Evolution offers marginal improvement (+0.9pp) under stable configuration. Structure prior (DAN templates) dominates with +19.7pp.}
\label{tab:ablation}
\small
\begin{tabular}{lccc}
\toprule
\textbf{Component} & \textbf{ASR (\%)} & \textbf{$\Delta$} & \textbf{Insight} \\
\midrule
PAIR baseline (no skills) & 55.5 & — & Reference point \\
\midrule
\textbf{Skill System} & & & \\
+ Cold Start + Retrieval + Maintenance & 78.2 & +22.7 & Primary contribution \\
+ Evolution (statistical) & 79.1 & +0.9 & Marginal refinement \\
\midrule
\textbf{Structure Prior} & & & \\
DAN templates (no evolution) & 98.8 & +19.7 & Dominant factor \\
DAN + Evolution & 99.7 & +0.9 & Ceiling reached \\
\bottomrule
\end{tabular}
\end{table}

% Key insight for paper narrative:
% - Skill System architecture is the main engine (+22.7pp)
% - Evolution is designed for refinement, not breakthrough (+0.9pp under optimal config)
% - Structure prior (DAN) determines the performance ceiling
%
% Data sources:
% - PAIR baseline: 55.5% (transfer experiments)
% - CS + Retrieval + Maintenance: 78.2% (ablation: single_call + trajectory + statistical, skip evolution)
% - CS + Evolution: 79.1% (layer1: single_call + trajectory + statistical, best config)
% - DAN only: 98.8% (layer3: full_evolve, 6 templates unchanged)
% - DAN + Evolution: 99.7% (layer4: medium + evo)